% =============================================================================
\chapter{Identificação e Delimitação do Plano}
% =============================================================================

Este documento apresenta o Plano de Estudos e Pesquisa vinculado ao mestrado
em Ciência da Computação no PPGC/UFRGS. A proposta investiga a integração entre
planejamento HTN e aprendizado por reforço multiobjetivo para a tomada de
decisão de NPCs em jogos digitais.

\section{Dados gerais}

\begin{description}
\item[Aluno:] Anderson G. A. C. Filho.
\item[Programa:] Mestrado em Ciência da Computação, PPGC/UFRGS.
\item[Orientador:] Prof. Dr. Lucas Nunes Alegre.
\item[Projeto institucional:] \textit{Efficient Reinforcement Learning}.
\item[Tema:] seleção online de métodos de decomposição HTN por MORL condicionado por preferências.
\item[Domínio de aplicação:] tomada de decisão hierárquica de NPCs em ambientes de jogo controlados.
\item[Estado inicial:] baseline HTN funcional integrado a um GridWorld com execução por ticks, sensores, navegação e replanejamento.
\end{description}

\section{Contexto de pesquisa}

NPCs são frequentemente controlados por técnicas simbólicas e baseadas em
regras, como máquinas de estados, Behavior Trees \cite{colledanchise2018}, GOAP
e HTN, amplamente discutidas em Game AI \cite{millington2019}. Essas técnicas
fornecem previsibilidade, consistência semântica e controle ao projetista, mas
normalmente dependem de prioridades fixas ou heurísticas construídas
manualmente. A hibridização entre GOAP e HTN também aparece como alternativa
para planejamento em narrativas interativas \cite{artar2019}.
Em ambientes dinâmicos, o número de condições necessárias para representar todas
as variações de contexto pode crescer rapidamente.

O aprendizado por reforço permite adaptação por experiência, porém o
aprendizado direto de ações de baixo nível amplia os espaços de estado e ação,
pode exigir grandes volumes de interação e reduz a transparência da decisão.
O MORL acrescenta a possibilidade de representar explicitamente objetivos
conflitantes, como sobrevivência, cumprimento da missão, proteção de aliados,
conservação de recursos, furtividade e velocidade.

A proposta combina essas propriedades: o HTN define o que é semanticamente
válido; uma fonte de preferências fornece $\mathbf{w}_t$ para expressar o
compromisso vigente; e o MORL escolhe qual método HTN válido deve decompor a
tarefa composta corrente.
A geração de preferências é comparada como componente auxiliar; codificadores
são implementações opcionais dessa fonte, e não uma camada de planejamento
independente.

\section{Problema de pesquisa}

Planejadores HTN tradicionais selecionam métodos por ordem fixa, custos
manuais ou heurísticas estáticas.
Embora previsíveis, esses mecanismos podem ser inadequados quando a melhor
decomposição depende de saúde, munição, condição de aliados, força inimiga,
perigos, objetivo da missão, personalidade ou eventos ocorridos durante a execução.

Uma política puramente aprendida, por outro lado, pode selecionar
comportamentos inválidos, incoerentes ou difíceis de explicar.
O problema investigado é, portanto, como introduzir adaptação multiobjetivo na
escolha de métodos HTN sem perder a validade semântica, a rastreabilidade e a auditabilidade
e o controle de designer fornecidos pelo planejamento simbólico.

\section{Questão principal}

\begin{quote}
O aprendizado por reforço multiobjetivo condicionado por preferências pode
melhorar a seleção online de métodos de decomposição HTN para NPCs, preservando
validade semântica, rastreabilidade, auditabilidade e desempenho de decisão em tempo real?
\end{quote}

\section{Questões secundárias}

\begin{enumerate}
\item A arquitetura produz melhores compromissos multiobjetivo que prioridades HTN estáticas?
\item Uma única política pode adaptar decisões a diferentes vetores $\mathbf{w}_t$ sem retreinamento?
\item Um grafo relacional ou um modelo profundo melhora a inferência de $\mathbf{w}_t$ em relação a regras explícitas?
\item A filtragem simbólica reduz a complexidade de aprendizado e impede métodos inválidos?
\item As decisões permanecem consistentes com as prioridades correntes do agente?
\item Qual é o custo computacional da consulta MORL durante a decomposição online?
\end{enumerate}

\section{Delimitação do escopo}

O foco está no nível hierárquico de decisão.
A ação da política MORL é um método de decomposição HTN, não um comando motor.
Movimento, mira, animação, desvio de obstáculos, pathfinding e execução individual
de ataques serão realizados por controladores determinísticos ou sistemas
tradicionais de jogos.

O escopo experimental inicial será restrito a um agente principal, um domínio
controlado, três a cinco tarefas compostas, duas a cinco alternativas por
tarefa e três a cinco objetivos.
Coordenação multiagente fica fora da contribuição principal.
Uma fonte relacional de preferências é a primeira alternativa estrutural; um
codificador profundo é condicional aos resultados, aos dados e ao orçamento
experimental disponíveis.
